% F3 -- From a raw day-file to training sequences, and the disjoint user split.
\begin{figure}[t]
\centering
\fitwidth{%
\begin{tikzpicture}[font=\small,
  stage/.style={draw=clrule,fill=clsoft,rounded corners=2pt,align=center,
                inner sep=5pt,text width=3.15cm,minimum height=1.5cm},
  raw/.style={draw=clrule,fill=white,rounded corners=2pt,align=left,
              inner sep=5pt,font=\scriptsize\ttfamily,text width=14.4cm},
  disk/.style={draw=clmodel,fill=clmodel!10,circle,align=center,
               inner sep=2pt,minimum size=2.05cm,font=\footnotesize},
  disktest/.style={draw=clink,fill=clink!14,circle,align=center,
               inner sep=2pt,minimum size=2.05cm,font=\footnotesize},
  ar/.style={-{Stealth[length=4pt]},clrule,thick},
]
\node[raw] (raw) {892739;22;M;36452;CETR;;;136;BKVDOU3;23822;BKVVIV2;24105;BKVSTN4;24154;\dots};
\node[font=\scriptsize,anchor=north,text=clrule] at ([xshift=-4.9cm,yshift=-0.34cm]raw.south)
  {$\underbrace{\hspace{4.6em}}$};
\node[font=\scriptsize,anchor=north,text=clrule] at ([xshift=-4.9cm,yshift=-0.62cm]raw.south)
  {8 metadata columns};
\node[font=\scriptsize,anchor=north,text=clrule] at ([xshift=1.9cm,yshift=-0.34cm]raw.south)
  {$\underbrace{\hspace{16em}}$};
\node[font=\scriptsize,anchor=north,text=clrule] at ([xshift=1.9cm,yshift=-0.62cm]raw.south)
  {(\cid{cell\_id}, timestamp) pairs --- the timestamp always \emph{follows} its cell};

\node[stage,below left=1.95cm and 3.4cm of raw.south] (s1)
  {\textbf{1. Keep what is a\\trajectory}\\[1pt]\footnotesize user id, record count,\\ the (cell, timestamp) pairs};
\node[stage,right=0.5cm of s1] (s2)
  {\textbf{2. Second $\rightarrow$ hour}\\[1pt]\footnotesize
   $\lfloor \mathit{ts}/3600 \rfloor$\\ one hour token per event};
\node[stage,right=0.5cm of s2] (s3)
  {\textbf{3. One line per user}\\[1pt]\footnotesize an ordered sequence of
   (\cid{cell\_id}, hour) events for one day};
\node[stage,right=0.5cm of s3] (s4)
  {\textbf{4. Split by user}\\[1pt]\footnotesize stratified on the number of
   records, \emph{before} any training};

\draw[ar] ([yshift=-0.95cm]raw.south) -- ++(0,-0.35) -| (s1.north);
\draw[ar] (s1) -- (s2); \draw[ar] (s2) -- (s3); \draw[ar] (s3) -- (s4);

\node[disk,below=1.25cm of s2] (tr) {\textbf{TRAIN}\\[2pt]{\large 400}\\[1pt]users the\\model studies};
\node[disktest,below=1.25cm of s4] (te) {\textbf{TEST}\\[2pt]{\large 100}\\[1pt]users it is\\graded on};
\node[font=\footnotesize,text=clloss,align=center] (neq) at ($(tr)!0.5!(te)$)
  {$\neq$\\[1pt]\textbf{0 user}\\ in common};

\draw[ar] (s4.south) -- ++(0,-0.4) -| (te.north);
\draw[ar] (s4.south) -- ++(0,-0.4) -| (tr.north);

\node[draw=clrule,fill=white,rounded corners=2pt,align=left,inner sep=6pt,
      text width=14.4cm,font=\footnotesize,below=0.75cm of neq]
  {\textbf{What the split guarantees, and what it does not.} The model never
   sees a single test user before being graded, so nothing it scores is
   memorised. But a user is \emph{one day}, and the day-files share no users at
   all: the model can only ever learn how a \emph{population} moves, never how
   one person moves. This single fact explains most of the ceilings reported in
   Chapters~\ref{ch:user} and~\ref{ch:scale}.};
\end{tikzpicture}}
\caption[From a raw day-file to training sequences]{How a raw daily export
becomes training material. The 400/100 figures are those of the
mobile-user day; the other tiers of Table~\ref{tab:datasets}
follow the same four stages with different counts.}
\label{fig:data-flow}
\end{figure}
